% Corps partagé par variants/answerspace-{none,inline}.tex : espace de réponse.
% Les marqueurs [A]…[I] permettent à check-answerspace.sh de mesurer les écarts.
\begin{exercise}
    Sans réglage, aucune question ne laisse de place.
    \begin{question}
        [A] Question sans espace.
    \end{question}
    \begin{question}
        [B] Question suivante.
    \end{question}
\end{exercise}

\setanswerspace{3cm}
\begin{exercise}
    Valeur par défaut du document : 3cm.
    \begin{question}
        [C] Question avec la valeur par défaut.
    \end{question}
    \begin{correction}
        Correction de C : elle remplace la place laissée dans le sujet.
    \end{correction}
    \begin{question}[space=0pt]
        [D] Question sans espace malgré la valeur par défaut.
    \end{question}
    \begin{question}[space=5cm]
        [E] Question avec sa propre valeur, 5cm.
    \end{question}
    \begin{question}
        [F] Question découpée :
        \begin{subquestion}
            [G] sous-question sans espace (pas de défaut pour les sous-questions) ;
        \end{subquestion}
        \begin{subquestion}[space=2cm]
            [H] sous-question avec 2cm.
        \end{subquestion}
    \end{question}
\end{exercise}

\begin{exercise}[answerspace=1cm]
    Valeur de l'exercice : 1cm.
    \begin{question}
        [I] Question avec la valeur de l'exercice.
    \end{question}
    \begin{question}
        [J] Dernière question.
    \end{question}
    \answerspace[2cm]
    [K] Après un espace isolé de 2cm.
\end{exercise}
\setanswerspace{0pt}
